\documentclass[letterpaper,11pt]{article}

\usepackage{fontawesome5}
\usepackage{latexsym}
\usepackage[empty]{fullpage}
\usepackage{titlesec}
\usepackage{marvosym}
\usepackage[usenames,dvipsnames]{color}
\usepackage{verbatim}
\usepackage{enumitem}

\usepackage[hidelinks]{hyperref}

\usepackage{fancyhdr}
\usepackage[english]{babel}
\usepackage{tabularx}
\input{glyphtounicode}


%----------FONT OPTIONS----------
% sans-serif
% \usepackage[sfdefault]{FiraSans}
% \usepackage[sfdefault]{roboto}
% \usepackage[sfdefault]{noto-sans}
\usepackage[default]{sourcesanspro}

% serif
% \usepackage{CormorantGaramond}
% \usepackage{charter}


\pagestyle{fancy}
\fancyhf{} % clear all header and footer fields
\fancyfoot{}
\renewcommand{\headrulewidth}{0pt}
\renewcommand{\footrulewidth}{0pt}

% Adjust margins
\addtolength{\oddsidemargin}{-0.5in}
\addtolength{\evensidemargin}{-0.5in}
\addtolength{\textwidth}{1.0in}
\addtolength{\topmargin}{-.5in}
\addtolength{\textheight}{1.0in}

\urlstyle{same}

\raggedbottom
\raggedright
\setlength{\tabcolsep}{0in}

\titleformat{\section}{
  \vspace{2pt}\scshape\raggedright\large
}{}{0em}{}[\color{black}\titlerule \vspace{-5pt}]

% Ensure that generate pdf is machine readable/ATS parsable
\pdfgentounicode=1

%-------------------------
% Custom commands
\newcommand{\resumeItem}[1]{
  \item\small{
    {#1 \vspace{0pt}}
  }
}

\newcommand{\resumeSubheading}[4]{
  \vspace{0pt}\item
    \begin{tabular*}{0.97\textwidth}[t]{l@{\extracolsep{\fill}}r}
      \textbf{#1} & #2 \\
      \textit{\small#3} & \textit{\small #4} \\
    \end{tabular*}\vspace{-7pt}
}

\newcommand{\resumeSubSubheading}[2]{
  \vspace{0pt}\item
    \begin{tabular*}{0.97\textwidth}{l@{\extracolsep{\fill}}r}
      \textit{\small#1} & \textit{\small #2} \\
    \end{tabular*}\vspace{-7pt}
}

\newcommand{\resumeProjectHeading}[2]{
    \item
    \begin{tabular*}{0.97\textwidth}{l@{\extracolsep{\fill}}r}
      \small#1 & #2 \\
    \end{tabular*}\vspace{-7pt}
}

\newcommand{\resumeSubItem}[1]{\resumeItem{#1}\vspace{-4pt}}

\renewcommand\labelitemii{$\vcenter{\hbox{\tiny$\bullet$}}$}


\newcommand{\resumeSubHeadingListStart}{\begin{itemize}[leftmargin=0.15in, label={}]}
\newcommand{\resumeSubHeadingListEnd}{\end{itemize}}
\newcommand{\resumeItemListStart}{\vspace{0pt} \begin{itemize}[itemsep=0pt, parsep=0pt]}
\newcommand{\resumeItemListEnd}{\end{itemize}\vspace{-5pt}}

%-------------------------------------------
%%%%%%  RESUME STARTS HERE  %%%%%%%%%%%%%%%%%%%%%%%%%%%%


\begin{document}

\begin{center}
    \textbf{\Huge \scshape Chenyu Gao} \\ \vspace{5pt}
    \small \faPhone\ +1 (667) 910-5300 $~|~$ 
    \faEnvelope\ \href{mailto:chenyu.gao@vanderbilt.edu}{\underline{chenyu.gao@vanderbilt.edu}} $~|~$ 
    \faGlobe\ \href{https://chenyugoal.github.io/}{\underline{chenyugoal.github.io}} $~|~$
    \faLinkedin\ \href{https://www.linkedin.com/in/chenyu-gao-4633b31a4/}{\underline{LinkedIn}} $~|~$
    \faGoogle\ \href{https://scholar.google.com/citations?hl=en&user=rKHNyWoAAAAJ&view_op=list_works&sortby=pubdate}{\underline{Google Scholar}}
\end{center}


\section{Education}
  \resumeSubHeadingListStart
    \resumeSubheading
      {Vanderbilt University}{Nashville, TN}
      {Ph.D. in Electrical \& Computer Engineering}{July 2022 -- Nov 2026 (expected)}
      \resumeItemListStart
      \resumeItem{\textbf{Research Focus:} Applying multimodal representation learning and generative models to solve challenges in medical image analysis and computer vision.}
      
      % \resumeItem{Research area: multimodal representation learning, computer vision, generative models, medical imaging}
      \resumeItemListEnd
    \resumeSubheading
      {Johns Hopkins University}{Baltimore, MD}
      {M.S. in Biomedical Engineering}{Aug 2020 -- May 2022}
    \resumeSubheading
      {Sun Yat-sen University}{Guangzhou, China}
      {B.S. in Biomedical Engineering}{Aug 2016 -- June 2020}
  \resumeSubHeadingListEnd


\section{Experience}
\resumeSubHeadingListStart
    \resumeSubheading{Data Science and Machine Learning Intern}{South San Francisco, CA}{insitro}{June 2025 -- Aug 2025}
    \resumeItemListStart
        \resumeItem{Designed and implemented a multimodal framework applying natural language processing (NLP) to encode DNA sequencing data, guiding self-supervised learning (SSL) on images to extract robust features for gene discovery.}
        \resumeItem{Collaborated with research scientists to design model architecture and partnered with software engineers to use and help debug infrastructure pipelines for large-scale computing.}
        \resumeItem{First-authored a technical paper, currently pending internal review for its potential to accelerate drug discovery by identifying new gene targets.}
    \resumeItemListEnd
    
    \resumeSubheading{Graduate Research Assistant}{Nashville, TN}{Vanderbilt University}{July 2022 -- Present}
    \resumeItemListStart
        \resumeItem{Engineered a cascaded diffusion model for a privacy red-team attack, reconstructing high-fidelity 3D facial geometry from defaced MRI data to quantify and expose critical re-identification risks. \href{https://chenyugoal.github.io/research_projects/2024_DREAM/}{[Blog]}}
        \resumeItem{Designed a first-of-its-kind system for brain age estimation from multi-modal MRI to enable early detection of neurodegenerative diseases, resulting in a provisional patent. \href{https://github.com/MASILab/BRAID}{[GitHub]}}
        \resumeItem{Architected a scalable data processing pipeline using Singularity, local and HPC to ingest, harmonize, and quality-assure 100,000 MRI scans from 40+ datasets for the world's largest diffusion and structural MRI database.}
        \resumeItem{Implemented a conditional GAN in PyTorch to synthetically extend the MRI field-of-view, rescuing previously incomplete scans and increasing effective dataset size by an estimated 10-15\%.}
    \resumeItemListEnd

    \resumeSubheading{Graduate Research Assistant}{Baltimore, MD}{Johns Hopkins University}{Dec 2020 -- May 2022}
    \resumeItemListStart
        \resumeItem{Adapted a lifelong learning algorithm from vision to speech and validated its omnidirectional knowledge transfer on a spoken digit benchmark, contributing to a publication in the top-tier journal \textit{IEEE TPAMI}. \href{https://github.com/neurodata/ProgLearn}{[GitHub]}}
        \resumeItem{Developed and benchmarked a suite of deep learning and classical computer vision algorithms for MR image analysis, establishing performance baselines that guided subsequent research on MRI defacing.}
    \resumeItemListEnd
\resumeSubHeadingListEnd

\section{Recent Publications \& Intellectual Property}
\resumeSubHeadingListStart
    \item \small \textbf{Note:} For a complete list, please see my \href{https://scholar.google.com/citations?hl=en&user=rKHNyWoAAAAJ&view_op=list_works&sortby=pubdate}{\underline{Google Scholar}} profile. Visit my \href{https://chenyugoal.github.io/}{\underline{website}} for showcases.
    \resumeItemListStart
        \resumeItem{C Gao, et al. ``Pitfalls of defacing whole-head MRI: re-identification risk with diffusion models and compromised research potential.'' \textit{Computers in Biology and Medicine}. 2025.}
        \resumeItem{C Gao, et al. ``Brain age identification from diffusion MRI synergistically predicts neurodegenerative disease.'' \textit{Imaging Neuroscience}. 2025.}
        \resumeItem{\textbf{U.S. Patent Pending:} System and Method of Brain Age Identification... (Non-provisional application filed Sept. 2025)}
    \resumeItemListEnd
\resumeSubHeadingListEnd


\section{Technical Skills}
 \begin{itemize}[leftmargin=0.15in, label={}, itemsep=0pt, topsep=1pt]
    \item \small
        \begin{tabularx}{0.97\textwidth}{@{} >{\bfseries}l @{\hspace{20pt}} X @{}}
            Languages & Python, Bash \\
            Frameworks & PyTorch, TensorFlow, Scikit-learn, Hugging Face, Pandas, NumPy, Polars \\
            DevOps \& Cloud & Git, Docker, Singularity, HPC/Slurm, AWS (S3, EC2), Weights \& Biases \\
        \end{tabularx}        
 \end{itemize}
 
 \end{document}
